\begin{abstract}
\noindent
Do paragliders move differently when they fly in a group? Group time differs from solo flight in
many ways, but so do the pilots who fly in groups, the places and hours where groups form, and the
moments pilots choose to fly together. We measure the motion by the local exponent $\Hloc(\tau)$ of
the time-averaged mean squared displacement and by statistics of the climbs and glides, and take
each of the other causes away with one test: contrasts within strata of site, day and hour and
within the same flight; solo flights reweighted to the zones and hours of the groups; the same
flight before, during and after a group. Over a quarter of an hour a pilot in a group progresses
straighter: $\Delta\Hloc(1000~\mathrm{s})=\DHlocThousandFlight$ within the same flight and
\DHlocThousandGroupSdh{} at the same site, day and hour, of which who flies in groups accounts for
\DHlocThousandSel. The zone and the hour remove at most \ZoneRedLocZoneHour\% of the gap, and on the
same flight the effect starts and stops with the group. Within one glide or climb nothing changes.
In a group pilots also climb more, end more searches in a climb, and glide faster and flatter.
Most of this is there with two pilots, and the pilot ahead finds no more thermals than the one
behind. The tests fix the pilot, the place, the hour and the moment, but not what brings pilots
together: the effect goes with flying together, and whether the group causes it stays open.
\end{abstract}
